% !TeX root = ../main.tex
%------------------------------------------------------------------------------
%  Ground Improvement, Pre-Support and Construction Hazards
%  Sources: Kirsch & Bell, Ground Improvement; EN 12715:2020 and EN 12716:2018
%  on grouting; Andersland & Ladanyi on frozen ground; ITA-AITES WG5 guidelines
%  and BS 6164:2019 on health and safety in tunnelling.
%------------------------------------------------------------------------------
\chapter{Ground Improvement, Pre-Support and Construction Hazards}
\label{ch:ground-improvement}

Ground improvement and hazard control are treated in one chapter because they
share a logic. Improvement modifies the ground so that the intended excavation
method will function; hazard control manages conditions the design did not
allow for. Both are selected from the geological model, both consume programme
as well as money, and both are substantially cheaper as planned
measures than as responses. Both are therefore treated here as design
decisions rather than construction contingencies: a treatment programme
devised after a machine has stopped is worth a fraction of the same programme
devised before the drive began, and a probe drilling policy adopted after an
inrush comes too late.

\section{Treatment objectives and grouting}
\label{sec:treatment-grouting}

\subsection{Objectives and technique selection}
\label{sec:objectives-selection}

Four objectives are distinguishable, and confusing them is the commonest cause
of an ineffective treatment. Improvement may aim to increase \emph{strength},
so that the face or the excavation stands; to increase \emph{stiffness}, so
that deformation is reduced; to reduce \emph{permeability}, so that inflow or
drawdown is controlled; or to \emph{compensate movement}, so that a structure
above is protected. A technique excellent for one may be useless for another,
so the objective is decided before the technique is specified.

Selection then follows the ground (Table~\ref{tab:technique-selection}).
Permeation grouting is bounded by grading,
which decides whether a grout can enter the pores at all (the groutability
criterion of Section~\ref{sec:grouting-methods}): cement suspensions permeate
coarse sands and gravels, microfine cements finer sands, and chemical solutions
finer still, while clays cannot be permeated. Permeability, the
hydraulic expression of the same pore structure, sets the inflow objective
and therefore the acceptance criterion of the treatment.
Jet grouting is largely independent of grading and works by erosion and
mixing. Freezing works in any saturated ground and is limited by flowing
groundwater, which advects heat into the freeze zone faster than the plant
can remove it~\citep{andersland2004frozen}.
Soil mixing suits soft soils and is defeated by obstructions.

\begin{table}[htbp]
\caption{Selection of the treatment techniques of this chapter: what each does, the ground it suits, and its principal limit.}
\label{tab:technique-selection}
\centering
\small
\begin{tabularx}{\textwidth}{@{}>{\RaggedRight\arraybackslash}p{0.16\textwidth} Y Y Y@{}}
\toprule
\textbf{Technique} & \textbf{Action and use} & \textbf{Ground it suits} & \textbf{Principal limit} \\
\midrule
Permeation grouting & Fills the pore space without disturbing the soil structure & Coarse sands and gravels with cement, finer sands with microfine cement, finer still with chemical solutions & Grading: reliable at a groutability ratio above about 24; clays cannot be permeated \\
Rock grouting & Curtain grouted ahead of the face for inflow control & Rock ahead of the face & Accepted on permeability in test holes, not on grout take \\
Jet grouting & Erodes the soil and mixes it with cement in place to form columns & Largely independent of grading & Column diameter predicted, not controlled; heave if the spoil return blocks \\
Compaction grouting & Stiff mortar displaces and densifies the soil; densification and controlled heave & Soils that drain during injection & Unreliable in saturated fine-grained soils \\
Compensation grouting & Injection between tunnel and structure in response to monitored settlement & Ground between a tunnel and the structure above & Needs access and real-time monitoring; no protection against a sudden loss \\
Ground freezing & Converts pore water to ice, giving strength and effectively zero permeability & Any saturated ground, including fine soils under high head & Seepage above roughly 1~m/day; frost heave and thaw settlement; continuous energy supply \\
Deep soil mixing & Blends binder into the soil in place, for cut-off walls and block treatment & Soft soils & Obstructions; variable achieved strength \\
Umbrella arch and spiles & Carry the ground ahead of the face & Low-cover, weak-ground and portal work; spiles in ravelling ground & Occupies the critical path \\
Dewatering & Lowers the water table so that excavation proceeds in drained ground & Permeable soils by wellpoints and deep wells; low-permeability soils by ejectors & Drawdown that the surroundings will tolerate \\
\bottomrule
\end{tabularx}
\end{table}

Cost and programme differ by an order of magnitude across the range, and the
cheapest technique that achieves the stated objective in the ground present is
the one to select. Ground treatment cannot be specified reliably without a
permeability distribution.

\subsection{Grouting methods}
\label{sec:grouting-methods}

\emph{Permeation grouting} fills the pore space without disturbing the soil
structure, and its effectiveness is governed by groutability, conventionally
the ratio of the soil $D_{15}$ to the grout $D_{85}$, with permeation reliable
above about 24 and not achievable below about
11~\citep{kirschbell2013ground} (Figure~\ref{fig:groutability}). Sleeve pipes (tubes à manchette) allow
repeated, controlled injection at selected horizons, which makes
permeation grouting a controllable process rather than a single irreversible
injection. Pressure
and volume are both monitored, and pressure rising faster than volume suggests
that the grout is fracturing the ground rather than permeating it.

\begin{figure}[htbp]
\centering
\datafigure[0.62\textwidth]{ch25/groutability}
\caption{Groutability ratio $N = D_{15}/D_{85}$, soil $D_{15}$ over grout $D_{85}$. Permeation is reliable below the solid line ($N > 24$) and not achievable above the dashed line ($N < 11$), and uncertain between them, after \citet{kirschbell2013ground}.}
\label{fig:groutability}
\end{figure}

\emph{Rock grouting} ahead of the face serves inflow control rather than
strength. A curtain of holes drilled outward from the face is grouted before
excavation advances into it, and acceptance is set on a permeability criterion
measured in test holes rather than on grout take~\citep{en12715}. Systematic probe drilling
and pre-excavation grouting are established practice in hard rock tunnelling,
where the grouting programme is specified to a residual inflow per unit length
of tunnel that the lining and drainage are then designed to accept.

\emph{Jet grouting} erodes the soil with a high-velocity fluid jet and mixes
it with cement in place, forming columns whose diameter depends on soil type,
energy and system: single fluid, double fluid with an air shroud, or triple
fluid with separate erosion and grouting~\citep{en12716} (Figure~\ref{fig:jet-grouting}). It is powerful but imprecise: column
diameter is predicted rather than controlled, and verification by coring is
necessary where the column geometry matters. Two hazards recur. Spoil return
must be maintained, because a blocked return pressurises the ground and causes
uncontrolled heave. Jet grouting near an existing structure can damage it
through the same mechanism.

\begin{figure}[htbp]
\centering
\datafigure[\textwidth]{ch25/jet-grouting}
\caption{Jet grouting systems (schematic). The rod is rotated and withdrawn while the jets erode and mix the soil, forming the column below the monitor, and spoil returns up the annulus. An air shroud extends the jet, and the triple system erodes with water above a separate grout jet, giving progressively larger columns.}
\label{fig:jet-grouting}
\end{figure}

\subsection{Displacement and compensation grouting}
\label{sec:displacement-compensation}

\emph{Compaction grouting} injects a stiff mortar that displaces and densifies
the surrounding soil rather than permeating it, and it is used for
densification and for controlled heave. Densification depends on drainage
during injection, so the method is unreliable in saturated fine-grained soils,
where injection raises pore pressure instead; surface heave is the variable
by which the process is controlled.

\emph{Compensation grouting} is the specialised application that protects
structures during tunnelling. Grout is injected from an array of sleeve pipes
installed between the tunnel and the structure above
(Figure~\ref{fig:compensation-grouting}), in response to
monitored settlement, so that ground loss is intercepted before it reaches the
surface. It runs in three phases: conditioning to pre-stress the soil and
establish the response, concurrent injection during tunnelling to compensate
in real time, and corrective injection afterwards to recover residual
settlement.

\begin{figure}[htbp]
\centering
\datafigure[\textwidth]{ch25/compensation-grouting}
\caption{Compensation grouting (schematic). Sleeve pipes installed from a shaft between the tunnel and the building inject grout in response to monitored settlement, so that the ground loss at the tunnel does not reach the foundation as the greenfield trough would.}
\label{fig:compensation-grouting}
\end{figure}

It is an observational technique in the strict sense of
Chapter~\ref{ch:monitoring} (injection is driven
by monitoring rather than by a predetermined schedule) and it requires
real-time instrumentation, an established injection response, and a decision
process able to act within hours. Its limits should be stated: it compensates
settlement, it cannot reverse damage already done, it cannot protect against a
sudden loss, and it requires physical access between tunnel and structure that
is not always available.

\section{Freezing, mixing and pre-support}
\label{sec:freezing-presupport}

\subsection{Freezing and soil mixing}
\label{sec:freezing-mixing}

\emph{Ground freezing} converts pore water to ice, producing a material with
substantial strength and effectively zero permeability, and it works in ground
where other methods do not: fine saturated soils under high head. Brine
systems circulate refrigerant through freeze pipes and take weeks to form the
wall; liquid nitrogen forms it in days at much higher operating cost, and is
used for short durations or as a rescue measure.

Three constraints govern. Flowing groundwater supplies heat by advection and
can prevent closure of the freeze wall (Figure~\ref{fig:freeze-wall}); a seepage velocity above
roughly 1\,m/day makes freezing difficult, and the gap opens where the flow is,
typically in the most permeable layer (Figure~\ref{fig:r3d-freeze-wall}). Frost heave occurs in frost-susceptible
silts as ice lenses form, and thaw settlement follows when the ground is
allowed to warm, so structures above must be assessed for both. Finally, the freeze
wall must be maintained until the permanent works can carry the load, which
means continuous energy supply and monitoring by temperature sensors
throughout the wall.

\begin{figure}[htbp]
\centering
\datafigure[\textwidth]{ch25/freeze-wall}
\caption{Plan of a ring of freeze pipes around an excavation (schematic). Frozen columns grow around each pipe until they merge into a closed wall. With seepage at velocity $v$ the upstream columns stay small and are displaced downstream, leaving a gap through which water enters.}
\label{fig:freeze-wall}
\end{figure}

\begin{figure}[htbp]
\centering
\renderfigure{freeze_wall}{%
  \rlab[south]{0.37,0.62}{0.14,0.92}{Frozen wall}
  \rlab[south]{0.47,0.80}{0.55,0.96}{Freeze pipes}
  \rlab[east]{0.40,0.42}{0.22,0.18}{Unclosed gap}
  \rlab[west]{0.75,0.48}{0.80,0.22}{Permeable layer}
}
\caption{A freeze wall around a shaft, 16 pipes on a 3~m ring, with the front half of the ground cut away (schematic; depths illustrative). Columns of 0.45 to 1.05~m radius merge into a closed wall (malachite). In the permeable layer (light grey) seepage holds the five upstream columns (carmine) at 0.32 to 0.48~m, leaving the wall open between them.}
\label{fig:r3d-freeze-wall}
\end{figure}

\emph{Deep soil mixing} blends binder into soft soil in place, by auger or by
cutter, producing columns or panels for cut-off walls and block treatment.
Achieved strength is variable and substantially lower than that of laboratory
mixes of the same proportions, and uniformity is the quality issue: the specification
should be a characteristic strength from core testing rather than a nominal
mix.

\subsection{Pre-support and face treatment}
\label{sec:presupport-face}

Pre-support carries the ground ahead of the face so that excavation advances
under protection rather than into unsupported ground.

An \emph{umbrella arch} of steel pipes installed sub-horizontally around the
crown ahead of the face, grouted in place, spans from the last supported
section to the ground ahead. It is installed in overlapping stages
(Figure~\ref{fig:presupport}), each
extending several rounds beyond the face, and it is the standard measure for
low-cover, weak-ground and portal work (Section~\ref{sec:portals}). \emph{Forepoling spiles} are the
lighter version (bars driven ahead of the crown) suited to short-term
support of ravelling ground.

\emph{Face reinforcement} by fibreglass dowels installed into the advance core
increases its strength and stiffness and reduces extrusion. Fibreglass is used
because it can be cut through by the excavation without damaging the tools, and
the dowels are installed in overlapping campaigns and progressively excavated
away. This is the core-control principle from the design philosophies of
Chapter~\ref{ch:construction}, and its cost is the cycle time consumed by each
installation campaign.

\begin{figure}[htbp]
\centering
\datafigure[\textwidth]{ch25/presupport}
\caption{Longitudinal section at the face (schematic). Each stage of umbrella arch is installed from the crown and overlaps the previous one ahead of the face; fibreglass dowels reinforce the advance core and are excavated away round by round.}
\label{fig:presupport}
\end{figure}

Pre-support converts a face-stability problem into a
structural one (Figure~\ref{fig:r3d-umbrella-arch}), and it occupies the critical path in doing so. That programme
cost is why pre-support is designed for the reaches that need it rather than
applied continuously.

\begin{figure}[htbp]
\centering
\renderfigure{umbrella_arch}{%
  \rlab[south]{0.30,0.63}{0.25,0.96}{Previous stage}
  \rlab[south]{0.515,0.64}{0.49,0.96}{Overlap}
  \rlab[south]{0.70,0.66}{0.72,0.96}{Current stage}
  \rlab[north]{0.49,0.45}{0.45,0.10}{Face}
  \rlab[north]{0.66,0.355}{0.72,0.12}{Face dowels}
}
\caption{Two stages of umbrella arch over the face, with half the tunnel and core cut away (schematic). Each stage is 15~m of grouted pipe (malachite) inclined outward at 7~per~cent, and successive stages overlap by 4~m ahead of the face. Fibreglass dowels (pale malachite) reinforce the core. Profile and spacings illustrative.}
\label{fig:r3d-umbrella-arch}
\end{figure}

\subsection{Dewatering and verification}
\label{sec:dewatering-verification}

\emph{Dewatering} lowers the water table so that excavation proceeds in
drained ground. Wellpoints suit shallow drawdown in permeable soils; deep
wells suit larger drawdown and greater depth; ejectors suit low-permeability
soils where yields are small. The usual constraint is the
drawdown that the surroundings will tolerate
(Chapter~\ref{ch:hydrogeology}). Recharge wells can limit off-site drawdown
where the constraint binds, at the cost of a system that must be balanced
continuously.

\emph{Verification} is the step most often curtailed, although it is the only
evidence that the treatment has met its objective. Coring gives direct evidence of
strength and continuity; permeability testing verifies a cut-off; geophysics
gives continuity between verification points; and trial panels before the
production programme establish achievable parameters in the actual ground.

Acceptance criteria must be measurable and set before the work begins. A
specification requiring the treated ground to be ``adequate'' cannot be tested,
and its practical effect is that the treatment is accepted because it was
carried out rather than because it worked.

\section{Hazard identification, water and gas}
\label{sec:water-gas-hazards}

\subsection{Hazard identification}
\label{sec:hazard-identification}

Hazard identification is a systematic pass along the alignment asking, for
each reach, what could go wrong and what the consequence would be. It separates
geotechnical hazards (instability, inrush, squeezing) from construction
hazards such as fire and plant, and from occupational hazards such as
atmosphere and confined space, because each is managed by a different
mechanism.

The output is a register keyed to chainage, and its value depends on being
maintained. A register compiled at design stage and never revisited becomes a
compliance document; one updated as the ground is revealed remains a working
tool and feeds the baseline comparison of
Chapter~\ref{ch:delivery}.

\subsection{Water inrush and cavities}
\label{sec:inrush-cavities}

Inrush arises where the excavation intersects a confined high-head feature:
a karst conduit, a fault-hosted aquifer, a buried valley, or old flooded
workings. Flow rates can overwhelm pumping capacity within minutes, and the
heading is lost.

Precursors exist and should be watched for: increasing seepage, a change in
water clarity or temperature, drilling water losses, and softening at the
face. Probe drilling ahead of the face is the standard defence and should be
adopted as a policy through reaches identified as at risk, with a defined
standoff maintained between the probe depth and the face so that a feature is
detected before it is opened. Where a probe strikes water under pressure, it
is bled off in a controlled way and the ground pre-grouted before excavation
advances.

Cavities have to be located, because their positions cannot be predicted. Surface and borehole
geophysics can identify large voids; look-ahead seismic and electromagnetic
methods from the face give shorter-range warning. All have limits, and none
reliably detects a small water-filled cavity in advance, which is why
karst alignments are probed systematically rather than relying on geophysics.

\subsection{Gas hazards}
\label{sec:gas-hazards}

Methane from coal measures and organic sediments is the principal explosion
risk; hydrogen sulfide and carbon monoxide are acutely toxic; carbon dioxide
and oxygen deficiency cause asphyxiation; radon is a long-term exposure
hazard. Gas migrates along faults and permeable horizons and accumulates in
the crown, in dead ends and behind lining.

Classification determines the regime. A tunnel assessed as potentially gassy
or gassy attracts requirements that affect every part of the works:
continuous fixed monitoring with automatic alarms, personal monitors,
intrinsically safe electrical equipment, control of ignition sources including
hot work and battery charging, and ventilation designed for dilution rather
than only for air quality~\citep{bs6164,ita2008health}. Retrofitting that regime after gas is encountered
is disruptive and expensive, which is why classification is made from the
geological assessment before construction rather than from measurements at
the face.

\section{Atmosphere, fire and emergency control}
\label{sec:atmosphere-fire-emergency}

\subsection{Ventilation, heat and dust}
\label{sec:ventilation-heat-dust}

Construction ventilation dilutes dust, exhaust and gas, clears blasting fumes,
and removes heat. The required air quantity is set by whichever of these
duties governs, and in long deep drives it is usually heat removal. Forcing systems deliver fresh air to the face
through ducting and suit most drives; exhausting systems draw contaminated air
away and suit dust and fume control; combined systems handle both. Duct
leakage is the practical limit on drive length: a long duct loses a
substantial fraction of its flow, and ventilation capacity frequently
determines how far a single-entry drive can extend before a second access is
required.

Heat enters the drive from the rock, at a rate set by the virgin rock
temperature, which rises with depth. It also enters from the plant: a large machine dissipates most of its installed power as heat.
Beyond the capacity of ventilation to remove it, mechanical cooling is
required, and its power demand is a significant construction-stage carbon term
(Chapter~\ref{ch:sustainability}).

Respirable crystalline silica is the dominant health hazard in rock
excavation, causing irreversible disease after long latency. Control is by
water suppression at the source, local extraction, enclosed cabins with
filtered air, and exposure monitoring with health surveillance. Exposure
standards have tightened and continue to tighten, so a design that meets
the current standard marginally may not meet the standard in force when the
work is built.

\subsection{Contamination and fire}
\label{sec:contamination-fire}

Contaminated ground and groundwater change how spoil is handled, where it can
go and what it costs, and the classification is made from investigation
sampling rather than discovered during excavation. Asbestos-bearing rock,
acid sulfate soils and industrial contamination each carry specific handling,
treatment and disposal requirements, and each can dominate the spoil logistics
that Chapter~\ref{ch:sustainability} treats as a carbon term.

Fire underground is exceptionally dangerous because escape is one-dimensional
and smoke travels the same route. The fire load is largely plant (hydraulic
oil, tyres, cabling, conveyor belting) and the controls are correspondingly
practical: limiting combustible materials, fire-resistant hydraulic fluids and
belting, detection and suppression on plant, hot-work permits, and refuge
chambers within a defined travel distance of the face. Self-rescuers are
carried by everyone underground.

\subsection{Access, egress and emergency response}
\label{sec:access-egress-emergency}

Much underground work is legally a confined space, requiring atmospheric
testing before entry, permit systems and standby arrangements. The dominant
constraint on a driven tunnel is that it has one way out until it holes
through, so the length of a single-entry drive is limited by escape time as
well as by ventilation, and intermediate access shafts are provided for
egress as much as for construction.

Emergency planning must be specific. It requires a response
plan for each credible scenario (fire, inrush, collapse, gas, medical),
communications that survive a power failure, rescue capability with the
training and equipment to enter an irrespirable atmosphere, and rehearsal.
Rehearsal is necessary because the recurring
lesson from documented incidents is that response was delayed by unfamiliarity
rather than by lack of provision.

\chapterreferences
